\documentclass[
    language=english,
    doctype=review-article,
    institution=none,
]{../../omnilatex}

\RequirePackage{config/document-settings}

\begin{document}
\maketitle

\begin{abstract}
Transformer-based large language models (LLMs) have fundamentally reshaped
the landscape of natural language processing (NLP) and are increasingly
influencing computer vision, robotics, and scientific discovery. This
comprehensive review examines the architectural innovations, training
methodologies, and application domains that have defined the transformer
era from 2017 through mid-2025. We organize the literature thematically
around six core topics: (1) the original transformer architecture and
its variants, (2) pretraining objectives and scaling laws, (3) alignment
and fine-tuning techniques, (4) efficiency and deployment optimizations,
(5) multimodal extensions, and (6) safety, ethics, and societal impact.
For each theme, we synthesize findings from over 400 papers, identify
key open problems, and outline future research directions. Our analysis
reveals that while scaling continues to yield performance improvements,
fundamental challenges in reasoning, factual grounding, and alignment
remain largely unsolved.
\end{abstract}

\section{Introduction}
\label{sec:introduction}

The introduction of the transformer architecture by Vaswani et al.\
\cite{vaswani2017attention} marked a paradigm shift in sequence modeling.
By replacing recurrent and convolutional mechanisms with self-attention,
transformers enabled unprecedented parallelism during training and achieved
state-of-the-art results across a wide range of NLP benchmarks.

The subsequent development of pretrained language models --- BERT \cite{devlin2018bert},
GPT-2 \cite{radford2019language}, GPT-3 \cite{brown2020language}, and their
successors --- demonstrated that scaling model parameters, data, and compute
could yield emergent capabilities not explicitly present in smaller models.
By 2025, LLMs with hundreds of billions of parameters are used in production
systems serving hundreds of millions of users.

This review aims to provide a structured, thematic synthesis of the
transformer and LLM literature. Rather than cataloging individual papers,
we focus on identifying trends, contradictions, and open questions across
the field. The review is intended for researchers and practitioners seeking
a comprehensive overview of the current state of the art.

\section{The Transformer Architecture}
\label{sec:architecture}

\subsection{Original Architecture}

The transformer \cite{vaswani2017attention} consists of an encoder-decoder
structure, each composed of stacked self-attention and feed-forward layers.
The key innovation is the scaled dot-product attention mechanism:

\begin{equation}
    \text{Attention}(Q, K, V) = \text{softmax}\left(\frac{QK^T}{\sqrt{d_k}}\right) V
    \label{eq:attention}
\end{equation}

where $Q$, $K$, and $V$ are the query, key, and value matrices, and $d_k$
is the key dimension. Multi-head attention computes this function in parallel
across $h$ heads, with each head operating on a $d_k/h$-dimensional subspace.

\subsection{Encoder-Only Models}

BERT \cite{devlin2018bert} popularized the encoder-only transformer for
representation learning. Pretrained with masked language modeling (MLM) and
next sentence prediction (NSP), BERT achieved breakthrough results on a
range of NLU benchmarks. Subsequent variants include RoBERTa \cite{liu2019roberta}
(improved training procedure), ALBERT \cite{lan2019albert} (parameter
reduction), and DeBERTa \cite{he2020deberta} (disentangled attention).

\subsection{Decoder-Only Models}

The GPT series \cite{radford2019language,brown2020language} demonstrated that
autoregressive language modeling, combined with scale, could produce models
with remarkable few-shot learning abilities. Decoder-only architectures have
become dominant for generative tasks due to their simplicity and scalability.

\begin{table}[htbp]
    \centering
    \caption{Comparison of major transformer variants.}
    \label{tab:variants}
    \begin{tabular}{@{}lcccc@{}}
        \toprule
        Model & Architecture & Parameters & Training Data & Key Innovation \\
        \midrule
        BERT & Encoder & 340M & 160 GB & MLM pretraining \\
        GPT-2 & Decoder & 1.5B & 40 GB & Zero-shot transfer \\
        GPT-3 & Decoder & 175B & 300B tokens & In-context learning \\
        PaLM & Decoder & 540B & 780B tokens & Pathways system \\
        LLaMA 2 & Decoder & 70B & 2T tokens & Open-weight model \\
        GPT-4 & MoE (est.) & $\sim$1.8T & $\sim$13T tokens & Multimodal \\
        \bottomrule
    \end{tabular}
\end{table}

\section{Pretraining and Scaling}
\label{sec:pretraining}

\subsection{Scaling Laws}

Kaplan et al. \cite{kaplan2020scaling} established that language model
performance follows power-law relationships with model size, dataset size,
and compute budget:

\begin{equation}
    L(N) \propto N^{-\alpha_N}, \quad L(D) \propto D^{-\alpha_D}, \quad L(C) \propto C^{-\alpha_C}
    \label{eq:scaling}
\end{equation}

where $L$ is the loss, $N$ is the number of parameters, $D$ is the dataset
size, and $C$ is the compute budget. The Chinchilla study \cite{hoffmann2022chinchilla}
revised these findings, arguing that current LLMs are significantly
undertrained relative to their parameter count.

\begin{table}[htbp]
    \centering
    \caption{Chinchilla-optimal training configurations.}
    \label{tab:chinchilla}
    \begin{tabular}{@{}lcc@{}}
        \toprule
        Model Parameters & Optimal Tokens & Compute (FLOPs) \\
        \midrule
        1B  & 20B   & $6 \times 10^{19}$ \\
        10B & 200B  & $6 \times 10^{21}$ \\
        70B & 1.4T  & $4.2 \times 10^{23}$ \\
        175B & 3.5T & $1.05 \times 10^{24}$ \\
        \bottomrule
    \end{tabular}
\end{table}

\subsection{Pretraining Objectives}

The dominant pretraining objectives include:

\begin{itemize}
    \item \textbf{Causal language modeling (CLM):} Predict the next token
          given all preceding tokens. Used by GPT, LLaMA, and most generative
          models.
    \item \textbf{Masked language modeling (MLM):} Predict randomly masked
          tokens. Used by BERT and encoder-only models.
    \item \textbf{Prefix language modeling:} Bidirectional attention over a
          prefix, autoregressive generation thereafter. Used by T5 \cite{raffel2020exploring}
          and U-PaLM.
    \item \textbf{Fill-in-the-middle (FIM):} Train models to infill text
          given both prefix and suffix. Critical for code generation.
\end{itemize}

\section{Alignment and Fine-Tuning}
\label{sec:alignment}

\subsection{Instruction Tuning}

Instruction tuning \cite{ouyang2022training} fine-tunes pretrained LLMs on
collections of (instruction, response) pairs, improving their ability to
follow natural language instructions. Key datasets include FLAN \cite{wei2021finetuned},
Self-Instruct \cite{wang2023self}, and Dolly.

\subsection{Reinforcement Learning from Human Feedback (RLHF)}

RLHF \cite{ouyang2022training} aligns LLMs with human preferences through:
\begin{enumerate}
    \item Training a reward model on human preference data.
    \item Optimizing the LLM policy using PPO (Proximal Policy Optimization)
          to maximize the learned reward.
\end{enumerate}

\begin{equation}
    R_{\text{RLHF}}(\theta) = \mathbb{E}_{x,y \sim \pi_\theta}\left[r_\phi(x, y)\right] - \beta \cdot \text{KL}\left[\pi_\theta \| \pi_{\text{ref}}\right]
    \label{eq:rlhf}
\end{equation}

where $r_\phi$ is the reward model, $\pi_\theta$ is the policy, $\pi_{\text{ref}}$
is the reference policy (initial SFT model), and $\beta$ controls the KL
penalty.

\subsection{Direct Preference Optimization (DPO)}

DPO \cite{rafailov2023direct} eliminates the need for a separate reward model
by directly optimizing the policy on preference data:

\begin{equation}
    \mathcal{L}_{\text{DPO}}(\theta) = -\mathbb{E}_{(x, y_w, y_l)}\left[\log \sigma\left(\beta \log \frac{\pi_\theta(y_w|x)}{\pi_{\text{ref}}(y_w|x)} - \beta \log \frac{\pi_\theta(y_l|x)}{\pi_{\text{ref}}(y_l|x)}\right)\right]
    \label{eq:dpo}
\end{equation}

where $y_w$ and $y_l$ are the preferred and dispreferred responses.

\begin{table}[htbp]
    \centering
    \caption{Comparison of alignment methods.}
    \label{tab:alignment}
    \begin{tabular}{@{}lcccc@{}}
        \toprule
        Method & Reward Model & Training Cost & Stability & Open Source \\
        \midrule
        RLHF (PPO) & Yes & High & Moderate & Limited \\
        DPO & No & Low & High & Extensive \\
        RLAIF & AI-generated & Medium & High & Growing \\
        Constitutional AI & No & Medium & High & Yes \\
        KTO & No & Low & High & Yes \\
        \bottomrule
    \end{tabular}
\end{table}

\section{Efficiency and Deployment}
\label{sec:efficiency}

\subsection{Model Compression}

\begin{table}[htbp]
    \centering
    \caption{Model compression techniques and their impact.}
    \label{tab:compression}
    \begin{tabular}{@{}lccc@{}}
        \toprule
        Technique & Size Reduction & Speedup & Quality Impact \\
        \midrule
        INT8 quantization & 2$\times$ & 1.5--2$\times$ & Minimal \\
        INT4 quantization & 4$\times$ & 2--3$\times$ & Small--Moderate \\
        GPTQ & 4$\times$ & 2--3$\times$ & Small \\
        LoRA (rank-16) & 10--20$\times$ & 1$\times$ & Small \\
        Pruning (50\%) & 2$\times$ & 1.5--2$\times$ & Moderate \\
        Knowledge distillation & Variable & Variable & Moderate \\
        \bottomrule
    \end{tabular}
\end{table}

\subsection{Inference Optimization}

Key inference optimization strategies include:
\begin{itemize}
    \item \textbf{KV-cache:} Avoids recomputing key/value tensors for
          previously processed tokens.
    \item \textbf{Speculative decoding:} Uses a smaller draft model to
          propose tokens, verified by the larger model.
    \item \textbf{Continuous batching:} Dynamically groups requests to
          maximize GPU utilization.
    \item \textbf{Tensor parallelism:} Distributes model layers across
          multiple GPUs.
    \item \textbf{Flash Attention:} IO-aware exact attention algorithm that
          reduces memory footprint and improves throughput \cite{dao2022flashattention}.
\end{itemize}

\section{Multimodal Extensions}
\label{sec:multimodal}

The transformer architecture has been successfully extended to vision
(ViT \cite{dosovitskiy2020image}), audio (Whisper \cite{radford2023robust}),
and combined modalities. Vision-language models (VLMs) such as GPT-4V,
Gemini, and LLaVA integrate visual understanding with language generation.

\begin{table}[htbp]
    \centering
    \caption{Major multimodal models and their capabilities.}
    \label{tab:multimodal}
    \begin{tabular}{@{}lccl@{}}
        \toprule
        Model & Modalities & Parameters & Key Capability \\
        \midrule
        GPT-4V & Text, Image & $\sim$1.8T (est.) & Visual reasoning \\
        Gemini 1.5 & Text, Image, Audio, Video & $\sim$1.5T (est.) & 1M context \\
        LLaVA 1.6 & Text, Image & 13B & Open VLM \\
        Whisper & Audio & 1.5B & Speech recognition \\
        SAM 2 & Image, Video & 64M & Segmentation \\
        \bottomrule
    \end{tabular}
\end{table}

\section{Safety, Ethics, and Societal Impact}
\label{sec:safety}

\subsection{Harms and Risks}

LLMs present several categories of harm:
\begin{itemize}
    \item \textbf{Hallucination:} Generating plausible but factually incorrect
          information with high confidence.
    \item \textbf{Bias and fairness:} Reproducing or amplifying societal biases
          present in training data.
    \item \textbf{Misinformation:} Enabling mass production of synthetic text
          for deceptive purposes.
    \item \textbf{Privacy:} Memorizing and leaking personally identifiable
          information from training data.
    \item \textbf{Dual use:} Facilitating cyberattacks, social engineering,
          or weapons design.
\end{itemize}

\subsection{Mitigation Strategies}

\begin{table}[htbp]
    \centering
    \caption{Safety mitigation approaches.}
    \label{tab:safety}
    \begin{tabular}{@{}lp{10cm}@{}}
        \toprule
        Approach & Description \\
        \midrule
        Constitutional AI & Uses a set of principles to guide AI self-critique and revision \\
        RLHF/RLAIF & Aligns model outputs with human or AI-generated preferences \\
        Red teaming & Adversarial testing to identify failure modes \\
        Watermarking & Embeds detectable patterns in generated text \\
        Content filtering & Input/output classifiers to block harmful content \\
        Differential privacy & Provides formal guarantees on training data privacy \\
        \bottomrule
    \end{tabular}
\end{table}

\section{Future Directions}
\label{sec:future}

We identify five key research directions:

\begin{enumerate}
    \item \textbf{Reasoning:} Current LLMs struggle with multi-step logical
          reasoning, mathematical proof, and causal inference. Approaches
          like chain-of-thought prompting \cite{wei2022chain} show promise
          but remain fragile.
    \item \textbf{Long-context modeling:} While context windows have grown
          to 1M+ tokens, effective utilization of long contexts remains
          challenging.
    \item \textbf{Grounded generation:} Connecting LLM outputs to verifiable
          sources (retrieval-augmented generation, tool use) is critical for
          factual reliability.
    \item \textbf{Efficient scaling:} Mixture-of-experts, sparse attention,
          and novel architectures may enable continued scaling with reduced
          compute.
    \item \textbf{Evaluation:} Current benchmarks are saturating; new
          evaluation frameworks are needed to assess reasoning, factuality,
          and safety simultaneously.
\end{enumerate}

\section{Conclusion}
\label{sec:conclusion}

The transformer architecture has enabled a generation of language models with
unprecedented capabilities. However, the field faces fundamental challenges
in reasoning, alignment, and safety that cannot be solved by scaling alone.
We anticipate that the next phase of progress will come from architectural
innovation, improved training methodologies, and interdisciplinary
collaboration with safety researchers and ethicists.

\bibliographystyle{plain}
\bibliography{references}

\end{document}
